\documentclass[11pt,a4paper]{article}


\usepackage{comment}
\usepackage[margin=1.9cm]{geometry}
\usepackage{fontspec}
\usepackage{polyglossia}
\setdefaultlanguage{greek}
\setotherlanguage{english}
\setmainfont{DejaVu Serif}[ItalicFont={DejaVu Serif},ItalicFeatures={FakeSlant=0.2},BoldItalicFont={DejaVu Serif Bold},BoldItalicFeatures={FakeSlant=0.2}]
\setmonofont{DejaVu Sans Mono}[Scale=0.88]
\newfontfamily\greekfont{DejaVu Serif}[ItalicFont={DejaVu Serif},ItalicFeatures={FakeSlant=0.2},BoldItalicFont={DejaVu Serif Bold},BoldItalicFeatures={FakeSlant=0.2}]
\newfontfamily\greekfonttt{DejaVu Sans Mono}[Scale=0.88]
\usepackage{amsmath,amssymb,amsthm}
\usepackage{enumitem}
\usepackage{fancyvrb}
\usepackage{fancyhdr}
\usepackage[unicode,hidelinks]{hyperref}
\hypersetup{pdftitle={Ασκήσεις 6}}
\newcommand{\horrule}[1]{\rule{\linewidth}{#1}}
\newcommand{\eps}{\varepsilon}

\theoremstyle{definition}
\newtheorem{question}{Άσκηση}
\newcommand{\HALT}{\operatorname{HALT}}
\setlist[enumerate]{label=(\alph*),leftmargin=*,itemsep=3pt,topsep=5pt}
\setlength{\parindent}{0pt}
\setlength{\parskip}{5pt}
\setlength{\emergencystretch}{3em}
\pagestyle{fancy}
\fancyhf{}
\fancyfoot[L]{\small ΗΥ280 (2026) - Τσουρακάκης}
\fancyfoot[R]{\small\thepage}
\renewcommand{\headrulewidth}{0pt}
\DefineVerbatimEnvironment{pseudocode}{Verbatim}{fontsize=\small,xleftmargin=1em}


\begin{document}
\begin{center}
\horrule{0.4pt}\\[0.15cm]
{\large\textbf{ΗΥ 280 \textemdash{} Θεωρία Υπολογισμού}}\\[0.1cm]
{Ασκήσεις 6 \textemdash{} Γλώσσες}\\[0.1cm]
{Χαράλαμπος Ε. Τσουρακάκης}\\
\horrule{0.4pt}
\end{center}

\begin{center}\fbox{\begin{minipage}{0.93\linewidth}
\textbf{Οδηγίες για τις ασκήσεις}
\begin{itemize}[leftmargin=*,itemsep=6pt]
\item \textbf{Ασκήσεις 1-4:} Αποτελούν μέρος του παραδοτέου.
\end{itemize}
\end{minipage}}\end{center}\medskip

\begin{question}[$\bigstar$ Πληθάριθμος της συνένωσης]
Για πεπερασμένες γλώσσες $A,B$ υπάρχουν ακριβώς $|A|\,|B|$ ζεύγη
$(x,y)\in A\times B$, και κάθε ζεύγος παράγει τη συμβολοσειρά $xy$.

\begin{enumerate}
\item Για $A=\{0,1\}$ και $B=\{1,2\}$ βρείτε το $AB$ και συγκρίνετε το
$|AB|$ με το $|A|\,|B|$.

\item Για $A=\{0,01\}$ και $B=\{1,11\}$ βρείτε το $AB$.
Ισχύει $|AB|=|A|\,|B|$;

\item Αποδείξτε ότι για πεπερασμένες γλώσσες $A,B$ ισχύει
\[
|AB|\leq |A|\,|B|.
\]
\end{enumerate}
\end{question}
 

\begin{question}[$\bigstar\bigstar$ Η γλώσσα $\{0,10\}^*$]
Έστω
\[
L=\{0,10\}^*
\]
και έστω $K$ η γλώσσα που περιέχει την $\eps$ και όλες τις μη κενές
δυαδικές συμβολοσειρές που

\begin{enumerate}
\item τελειώνουν σε $0$, και
\item δεν περιέχουν την υποσυμβολοσειρά $11$.
\end{enumerate}

\begin{enumerate}
\item Γράψτε πέντε συμβολοσειρές του $L$ και επαληθεύστε ότι ανήκουν στο
$K$.

\item Αποδείξτε ότι $L\subseteq K.$

\item Αποδείξτε ότι $K\subseteq L.$


\item Συμπεράνετε ότι $L=K$
\end{enumerate}
\end{question}
 


\begin{question}[$\bigstar\bigstar$ Συμπλήρωμα συμβολοσειράς]
Για κάθε συμβολοσειρά $w\in\{0,1\}^*$, ορίζουμε το \emph{συμπλήρωμα}
$w^c$ αντικαθιστώντας κάθε $0$ με $1$ και κάθε $1$ με $0$.

Ισοδύναμα, η συνάρτηση συμπληρώματος ορίζεται αναδρομικά ως εξής:
\[
w^c=
\begin{cases}
\eps, & w=\eps,\\
1x^c, & w=0x,\\
0x^c, & w=1x.
\end{cases}
\]

Αποδείξτε ότι:

\begin{enumerate}
\item Για κάθε $w\in\{0,1\}^*$,
\[
|w|=|w^c|.
\]

\item Για κάθε $x,y\in\{0,1\}^*$,
\[
(xy)^c=x^cy^c.
\]

\item Για κάθε $w\in\{0,1\}^*$,
\[
\#(1,w)=\#(0,w^c).
\]

\item Για κάθε $w\in\{0,1\}^*$,
\[
(w^R)^c=(w^c)^R.
\]

\item Για κάθε $w\in\{0,1\}^*$ και κάθε $n\in\mathbb N$,
\[
(w^n)^c=(w^c)^n.
\]
\end{enumerate}
\end{question}
 

\begin{question}[$\bigstar\bigstar$ Αντιστροφή γλώσσας]
Υπενθυμίζουμε ότι η \emph{αντιστροφή} $w^R$ μιας συμβολοσειράς $w$
ορίζεται αναδρομικά ως εξής:
\[
w^R=
\begin{cases}
\eps, & w=\eps,\\
x^R a, & w=ax,
\end{cases}
\]
όπου $a\in\Sigma$ και $x\in\Sigma^*$.

Η \emph{αντιστροφή} μιας γλώσσας $L$ ορίζεται ως
\[
L^R=\{w^R\mid w\in L\}.
\]

\begin{enumerate}
\item[] Αποδείξτε ότι για κάθε γλώσσα $L$,
$ 
(L^*)^R=(L^R)^*.$

%\item Αποδείξτε ότι αν η $L$ είναι κανονική γλώσσα, τότε και η $L^R$
%είναι κανονική.
\end{enumerate}

Μπορείτε να χρησιμοποιήσετε χωρίς απόδειξη τις ακόλουθες ιδιότητες:
\[
L^*L^*=L^*, (w^R)^R=w,(xy)^R=y^Rx^R.
\]
\end{question}
 


\end{document}
